\documentclass[11pt]{article}

\usepackage[margin=1in]{geometry}
\usepackage{amsmath,amssymb,amsthm,mathtools}
\usepackage{graphicx}
\usepackage[hidelinks]{hyperref}

\newtheorem{theorem}{Theorem}
\newtheorem{lemma}{Lemma}

\DeclareMathOperator{\id}{id}
\newcommand{\BH}{\mathcal B(\mathcal H)}
\newcommand{\HH}{\mathcal H}
\newcommand{\one}{\mathbf 1}

\title{Finite-dimensional bistochastic quantum random variables are classical}
\author{Candidate partial result for arXiv:2005.00005}
\date{2026-06-21}

\begin{document}
\maketitle

\section*{Status}

This packet records a candidate partial result, likely valid.  It resolves the
finite-dimensional Hilbert-space case of a question raised by Plosker and
Ramsey in \cite{PR}.  It is not claimed to settle the full separable
Hilbert-space formulation.

\section*{Source question}

The source paper is Sergei Plosker and Brian Ramsey, ``Bistochastic operators
and quantum random variables'', arXiv:2005.00005 \cite{PR}.  In Section 5, after
Definition 5.1 and Theorem 5.5, the authors show that every classical
bistochastic operator on \(L^1(X,\mu)\) extends entrywise to a bistochastic
operator on \(L^1_{\HH}(X,\mu I_{\HH})\).  Immediately afterwards they write:

\begin{center}
\includegraphics[width=0.94\linewidth]{figures/open_problem_crop.png}
\end{center}

In source-line form, the key sentence is line 1088 of the parsed source:
\[
\text{``We have no example of a bistochastic operator on }
L^1_{\HH}(X,\mu I_{\HH})\text{ that does not arise in this way.''}
\]

\section*{Definitions}

Let \((X,\Sigma,\mu)\) be a finite measure space and let
\(\nu=\mu I_{\HH}\).  We use Plosker--Ramsey's definition: a linear operator
\(B\) on \(L^1_{\HH}(X,\nu)\) is \emph{bistochastic} if
\[
 B\geq 0,\qquad
 \int_X Bf\,d\mu=\int_X f\,d\mu\quad(f\in L^1_{\HH}(X,\nu)),
 \qquad
 B I_{\HH}= I_{\HH}.
\]
The source paper proves that a classical bistochastic operator \(S\) on
\(L^1(X,\mu)\) gives such a \(B\) by
\[
       B(fA)=S(f)A,\qquad f\in L^1(X,\mu),\ A\in \mathcal B(\HH).
\]

\section*{Idea of the proof}

The obstruction to exotic examples is already present on indicator functions.
Given two measurable sets \(E,F\), integrate \(B(\chi_E A)\) over \(F\).  This
produces a positive linear map \(\Theta_{F,E}:M_d\to M_d\).  Since \(B\)
preserves the full integral, \(\Theta_{X,E}=\mu(E)\id\), so
\(\Theta_{F,E}\) lies in the positive order interval
\([0,\mu(E)\id]\).

In finite dimensions that order interval is one-dimensional: every positive
map \(M_d\to M_d\) dominated by a scalar multiple of the identity map is itself
a scalar multiple of the identity.  Rank-one projections and the two-by-two
matrix-unit corners give an elementary proof.  Therefore every
\(\Theta_{F,E}\) is scalar, and the scalar coefficients form the transition
measures of an ordinary Markov/bistochastic operator \(S\) on \(L^1(X,\mu)\).
The identity \(B(\chi_E A)=S(\chi_E)A\) then extends from indicators to simple
functions and, by density in finite dimension, to all of
\(L^1_{\mathbb C^d}(X,\mu I_d)\).

\begin{lemma}[Order interval of the identity on matrices]\label{lem:order}
Let \(d<\infty\), and let \(\Phi:M_d(\mathbb C)\to M_d(\mathbb C)\) be a
positive linear map.  If \(c\,\id-\Phi\) is positive for some \(c\geq 0\), then
\(\Phi=\lambda\,\id\) for some \(0\leq \lambda\leq c\).
\end{lemma}

\begin{proof}
If \(c=0\) the claim is immediate, so scale to \(c=1\).  Let
\((e_i)_{i=1}^d\) be the standard basis and let \(E_{ij}\) be the matrix units.
For each \(i\), the rank-one projection \(E_{ii}\) satisfies
\[
        0\leq \Phi(E_{ii})\leq E_{ii}.
\]
Hence \(\Phi(E_{ii})=\lambda_iE_{ii}\) for some \(0\leq\lambda_i\leq 1\).

Fix \(i\neq j\).  Put
\[
 p_{\pm}=\frac12(E_{ii}+E_{jj}\pm E_{ij}\pm E_{ji}),
\]
the projections onto the lines spanned by \(e_i\pm e_j\).  Since
\(0\leq\Phi(p_{\pm})\leq p_{\pm}\), there are scalars
\(\alpha_{\pm}\in[0,1]\) such that \(\Phi(p_{\pm})=\alpha_{\pm}p_{\pm}\).
Adding the two identities and using \(p_++p_-=E_{ii}+E_{jj}\) gives
\[
 \lambda_iE_{ii}+\lambda_jE_{jj}
 =
 \frac{\alpha_++\alpha_-}{2}(E_{ii}+E_{jj})
 +\frac{\alpha_+-\alpha_-}{2}(E_{ij}+E_{ji}).
\]
Therefore \(\alpha_+=\alpha_-\) and \(\lambda_i=\lambda_j\).  Since \(i,j\)
were arbitrary, all \(\lambda_i\)'s are a common scalar \(\lambda\).

The same identity with \(p_+\) now yields
\[
        \Phi(E_{ij}+E_{ji})=\lambda(E_{ij}+E_{ji}).
\]
Using the rank-one projections onto \(e_i\pm i e_j\) gives
\[
        \Phi(iE_{ij}-iE_{ji})=\lambda(iE_{ij}-iE_{ji}).
\]
Combining the last two displayed identities gives
\(\Phi(E_{ij})=\lambda E_{ij}\) and \(\Phi(E_{ji})=\lambda E_{ji}\).  Thus
\(\Phi\) is \(\lambda\id\) on all matrix units.
\end{proof}

\begin{theorem}\label{thm:finite}
Let \((X,\Sigma,\mu)\) be a finite measure space and let
\(\HH=\mathbb C^d\).  If \(B\) is a bistochastic operator on
\(L^1_{\HH}(X,\mu I_d)\) in the sense of \cite{PR}, then there is a unique
classical bistochastic operator \(S\) on \(L^1(X,\mu)\) such that
\[
        B(fA)=S(f)A
        \qquad
        (f\in L^1(X,\mu),\ A\in M_d(\mathbb C)).
\]
Consequently, for finite-dimensional \(\HH\), every bistochastic operator on
\(L^1_{\HH}(X,\mu I_{\HH})\) is one of the entrywise classical extensions
constructed in \cite[Theorem 5.5]{PR}.
\end{theorem}

\begin{proof}
Let \(E,F\in\Sigma\).  For \(A\in M_d\), define
\[
        \Theta_{F,E}(A)=\int_F B(\chi_E A)\,d\mu .
\]
If \(A\geq 0\), then \(\chi_EA\geq0\), so \(B(\chi_EA)\geq0\) and hence
\(\Theta_{F,E}(A)\geq0\).  Thus \(\Theta_{F,E}\) is positive.  Since \(B\)
preserves the integral,
\[
        \Theta_{X,E}(A)
        =\int_X B(\chi_EA)\,d\mu
        =\int_X \chi_E A\,d\mu
        =\mu(E)A .
\]
Also
\[
        \mu(E)A-\Theta_{F,E}(A)=\Theta_{X\setminus F,E}(A)\geq0
        \qquad(A\geq0).
\]
By Lemma~\ref{lem:order}, there is a scalar \(a_E(F)\in[0,\mu(E)]\) such that
\[
        \Theta_{F,E}=a_E(F)\id_{M_d}.
\]

For fixed \(E\), the set function \(F\mapsto a_E(F)\) is a finite positive
measure: countable additivity follows from countable additivity of the
operator integral applied, for instance, to \(A=I_d\).  Moreover
\(a_E(F)=0\) whenever \(\mu(F)=0\).  Let
\[
        h_E=\frac{d a_E}{d\mu}\in L^1(X,\mu)_+
\]
be the Radon--Nikodym density.  Then for all \(F\in\Sigma\) and \(A\in M_d\),
\[
        \int_F B(\chi_EA)\,d\mu
        =a_E(F)A
        =
        \int_F h_EA\,d\mu .
\]
Since \(M_d\) is finite-dimensional, equality of all integrals over measurable
sets is equality almost everywhere.  Thus
\[
        B(\chi_EA)=h_EA\quad\text{a.e.}
\]

Define \(S\chi_E=h_E\), and extend linearly to scalar simple functions.  The
previous paragraph gives
\[
        B(sA)=S(s)A
\]
for every scalar simple function \(s\) and every \(A\in M_d\).  Positivity of
\(B\) gives positivity of \(S\) on simple functions.  Taking \(A=I_d\) and
using \(B I_d=I_d\) gives \(S1=1\).  Taking \(F=X\) above gives
\(\int_X S\chi_E\,d\mu=\mu(E)\), hence \(\int_X Ss\,d\mu=\int_Xs\,d\mu\) for
all simple \(s\).  Consequently \(S\) is contractive on scalar simple
functions in the compatible \(L^1\)-norm
\(\int(|\operatorname{Re}s|+|\operatorname{Im}s|)\,d\mu\), and therefore
extends uniquely to a positive integral-preserving operator on \(L^1(X,\mu)\).
The identities \(S1=1\) and \(\int Ss\,d\mu=\int s\,d\mu\) pass to the
extension, so \(S\) is classically bistochastic.

Finally let \(f\in L^1(X,\mu)\) and \(A\in M_d\).  Choose scalar simple
functions \(s_n\to f\) in \(L^1(X,\mu)\).  The source paper proves boundedness
of multiplication by fixed matrices on \(L^1_{\HH}(X,\mu I_{\HH})\)
\cite[Proposition 4.4]{PR}, and \(B\) is contractive
\cite[Proposition 5.4]{PR}.  Hence
\[
        B(s_nA)\to B(fA),\qquad S(s_n)A\to S(f)A
\]
in \(L^1_{\HH}(X,\mu I_d)\), and the identity passes to the limit:
\[
        B(fA)=S(f)A.
\]
For a general \(F\in L^1_{\mathbb C^d}(X,\mu I_d)\), write
\(F=\sum_{i,j=1}^d f_{ij}E_{ij}\) with scalar \(f_{ij}\in L^1(X,\mu)\); this is
the finite-dimensional instance of the coordinate description used in
\cite[Theorem 5.5]{PR}.  Then
\[
        B(F)=\sum_{i,j=1}^d S(f_{ij})E_{ij}.
\]
Thus \(B\) is exactly the entrywise extension of \(S\).  Uniqueness of \(S\)
is immediate from \(B(fI_d)=S(f)I_d\).
\end{proof}

\section*{Limitations}

The source sentence is stated for separable Hilbert spaces.  The theorem above
uses finite dimensionality twice: in the matrix order-interval lemma and in the
coordinate/density step extending from simple matrix-valued functions to the
whole \(L^1_{\mathbb C^d}\) space.  Thus this packet gives a full no-exotic
answer for finite-dimensional quantum systems, but leaves the infinite-
dimensional case open.

\section*{Verification and novelty check}

No computational verification is required.  The proof reduces the problem to
Lemma~\ref{lem:order}, which is checked directly on matrix units, and to the
standard Radon--Nikodym argument for scalar transition measures.

Local run indexes were searched for arXiv id \(2005.00005\), the title
keywords, and terms including ``bistochastic'', ``quantum random variables'',
``majorization'', and ``operator algebra''.  No existing packet or attempt
covered this result.  A bounded web search on 2026-06-21 for exact phrases
around ``We have no example of a bistochastic operator'' and
``\(L^1_{\HH}(X,\mu I_{\HH})\) bistochastic'' found no duplicate answer.

Human review should focus on the finite-dimensional order-interval lemma and
on the passage from the scalar measures \(a_E\) to the classical operator
\(S\).

\begin{thebibliography}{9}
\bibitem{PR}
S. Plosker and B. Ramsey,
\emph{Bistochastic operators and quantum random variables},
arXiv:2005.00005, 2020.
\end{thebibliography}

\end{document}
